% Tapis de Quillen, with Quillen's offprint and preprint (cote 162-5, typed part): archive cote 162-5, site page 53 (PDF page 54).
% offprint. Draft compared: none.
% AI-transcribed by Claude Opus 5.5 (claude-opus-5-5, adaptive thinking, effort high), Message Batches API, batch t162_full (msgbatch_01UAh3zyWPd454AeDkhnmQ33), 2026-11-13; not yet checked by eye.
\scanpage{162-5/53}
\mspage{240} D.G. Quillen:

and where the second may be derived similarly (see also [2], p. 153). Thus
\[ \|u\|^2_{q*} = \sum_{|\alpha|=q} (\alpha!)^{-1}\, \|z^\alpha u\|^2_0 = \sum_{|\alpha|=q} (\alpha!)^{-1} (u, D^\alpha z^\alpha u)_0 \]
\[ = \sum_{|\alpha|=q} \sum_{\gamma} \alpha!\,(\gamma!\,(\alpha-\gamma)!^2)^{-1} (u, z^{\alpha-\gamma} D^{\alpha-\gamma} u)_0 \]
\[ = \sum_{|\alpha|=q} \sum_{\gamma} C_\gamma{}^\alpha ((\alpha-\gamma)!)^{-1}\, \|D^{\alpha-\gamma} u\|^2_0 \]
from which the lemma immediately follows.

Proof of the Theorem. The implication (ii)$\Rightarrow$(i) is clear, so assume (i)
holds. If $u \in S_{m+N}\, T \otimes V$ then by Lemma 1, (2), and (5) we have with
indices running like $1 \leqq i,\ j \leqq r$, $|\alpha| = |\beta| = m$
\[ E^N(u,u) = \sum a_{\alpha i, \beta j} (D^\alpha u_i, D^\beta u_j)_0 \]
\[ = \sum a_{\alpha i, \beta j} (u_i, z^\alpha D^\beta u_j)_0 \]
\[ = \int \sum (-1)^{|\gamma|}\, a_{\alpha i, \beta j}\, \alpha!\, \beta!\, (\gamma!\, (\beta-\gamma)!\, (\alpha-\gamma)!)^{-1}\, z^{\beta-\gamma} u_i (z^{\alpha-\gamma} u_j)^*\, dG \]
\[ = \int [E^\wedge (z^*, u) + Q(z,u)]\, dG \]
where $Q(z,u)$ is of degree $<m$ in $z_i z_j^*$ and of degree 1 in $u_i u_j^*$. Thus
for some constant $C$
\[ |Q(z,u)| \leqq C \sum_{i<m} (i!)^{-1}\, |z|^{2i}\, |u|^2, \]
while by hypothesis on $E$ there is an $\varepsilon > 0$ such that
\[ |E^\wedge (z^*, u)| \geqq \varepsilon (m!)^{-1}\, |z|^{2m}\, |u|^2. \]
Therefore
\[ E^N(u,u) \geqq \int (\varepsilon (m!)^{-1}\, |z|^{2m} - C \sum_{i<m} (i!)^{-1}\, |z|^{2i})\, |u|^2\, dG \]
\[ = \varepsilon\, \|u\|^2_{m*} - C \sum_{i<m} \|u\|^2_{i*} \]
\[ \geqq \varepsilon\, \|u\|^2_m - C' \sum_{i<m} \|u\|^2_i \]
using Lemma 3. By Lemma 2 we have $E^N(u,u) \geqq P(N)\, \|u\|^2_0$ where
\[ P(N) = \varepsilon\, C_m{}^{m+N} - C' \sum_{i<m} C_i{}^{i+N} \]
is a polynomial of degree $m$ in $N$ with positive leading coefficient. Thus
for $N$ sufficiently large $P(N) > 0$ so $E^N$ is positive definite and the theorem
is proved.

We shall now give some applications of the theorem when $V = \mathbf{C}$
\margin{\unsure{$\backslash$}} and $T = \mathbf{C}^n$.

\mspage{\hand{53}}
